\section{Ядро Дирихле}\label{sec-fourier-dirichlet}

Пусть \(f\) абсолютно интегрируема на \([-\pi,\pi]\). Обозначим частичную сумму её тригонометрического ряда Фурье через
\begin{equation}
T_n(x)=\frac{a_0}{2}+\sum_{k=1}^n\bigl(a_k\cos kx+b_k\sin kx\bigr),
\qquad n\ge0.
\label{eq-fourier-partial-sum}
\end{equation}

\begin{definition}[Ядро Дирихле]\label{def-dirichlet-kernel}

Функция
\begin{equation}
D_n(u)=\frac12+\sum_{k=1}^n\cos ku
\label{eq-dirichlet-kernel}
\end{equation}
называется ядром Дирихле. При \(n=0\) сумма пуста и \(D_0(u)=1/2\).

\end{definition}

Подставим коэффициенты \eqref{eq-trigonometric-coefficients} в \eqref{eq-fourier-partial-sum}. По линейности интеграла и формуле \(\cos kx\cos kt+\sin kx\sin kt=\cos k(x-t)\) имеем
\[
\begin{aligned}
T_n(x)
&=\frac1\pi\int_{-\pi}^{\pi}f(t)
\left(\frac12+\sum_{k=1}^n\bigl(\cos kx\cos kt+\sin kx\sin kt\bigr)\right)\,dt\\
&=\frac1\pi\int_{-\pi}^{\pi}f(t)
\left(\frac12+\sum_{k=1}^n\cos k(x-t)\right)\,dt.
\end{aligned}
\]
Таким образом,
\begin{equation}
T_n(x)=\frac1\pi\int_{-\pi}^{\pi}f(t)D_n(x-t)\,dt.
\label{eq-fourier-kernel-integral}
\end{equation}

\begin{lemma}[Явная формула и свойства ядра]\label{lem-dirichlet-properties}

Ядро \(D_n\) непрерывно, чётно и \(2\pi\)-периодично. При \(u\notin2\pi\mathbb Z\)
\begin{equation}
D_n(u)=\frac{\sin\bigl((n+\tfrac12)u\bigr)}{2\sin(u/2)}.
\label{eq-dirichlet-sine-form}
\end{equation}
В точках \(u=2\pi m\), \(m\in\mathbb Z\), значение равно \(n+1/2\). Кроме того,
\begin{equation}
\frac1\pi\int_{-\pi}^{\pi}D_n(u)\,du=1,
\qquad
\frac2\pi\int_0^\pi D_n(u)\,du=1.
\label{eq-dirichlet-normalization}
\end{equation}

\end{lemma}

\begin{proof}

Непрерывность, чётность и периодичность следуют из определения. Для явной формулы используем тождество
\[
2\sin(u/2)\cos ku
=\sin\bigl((k+\tfrac12)u\bigr)-\sin\bigl((k-\tfrac12)u\bigr).
\]
После суммирования по \(k=1,\ldots,n\) члены сокращаются:
\[
\begin{aligned}
2\sin(u/2)D_n(u)
&=\sin(u/2)+\sum_{k=1}^n
\left(\sin\bigl((k+\tfrac12)u\bigr)-\sin\bigl((k-\tfrac12)u\bigr)\right)\\
&=\sin\bigl((n+\tfrac12)u\bigr).
\end{aligned}
\]
При \(u\notin2\pi\mathbb Z\) делим на \(2\sin(u/2)\); при \(u=2\pi m\) используем определение \eqref{eq-dirichlet-kernel}. Интеграл каждого \(\cos ku\), \(k\ge1\), на \([-\pi,\pi]\) равен нулю, а интеграл постоянного члена \(1/2\) равен \(\pi\). Второе равенство \eqref{eq-dirichlet-normalization} следует из чётности.

\end{proof}

\section{Принцип локализации}\label{sec-fourier-localization}

Далее \(f\) --- \(2\pi\)-периодическая функция, абсолютно интегрируемая на одном периоде. При задании функции на \([-\pi,\pi]\) используется её периодическое продолжение; значения в отдельных точках не влияют на коэффициенты и интегралы.

\begin{lemma}[Симметричная запись частичной суммы]\label{lem-fourier-symmetric-integral}

Для любого \(x\in\mathbb R\) справедливо равенство
\begin{equation}
T_n(x)=\frac1\pi\int_0^\pi\bigl(f(x-t)+f(x+t)\bigr)D_n(t)\,dt.
\label{eq-fourier-symmetric-integral}
\end{equation}

\end{lemma}

\begin{proof}

В \eqref{eq-fourier-kernel-integral} сделаем замену \(u=x-t\):
\[
T_n(x)=\frac1\pi\int_{x-\pi}^{x+\pi}f(x-u)D_n(u)\,du.
\]
Подынтегральная функция \(2\pi\)-периодична, поэтому интеграл по любому отрезку длины \(2\pi\) одинаков. Переносим пределы на \([-\pi,\pi]\), разбиваем интеграл в нуле и в отрицательной части заменяем \(u\) на \(-u\). По чётности \(D_n\) получаем
\[
\begin{aligned}
T_n(x)&=\frac1\pi\int_{-\pi}^{\pi}f(x-u)D_n(u)\,du\\
&=\frac1\pi\int_0^\pi\bigl(f(x-u)+f(x+u)\bigr)D_n(u)\,du.
\end{aligned}
\]

\end{proof}

\begin{theorem}[Принцип локализации Римана]\label{thm-fourier-localization}

Для любого \(x_0\in\mathbb R\) и любого фиксированного \(\delta\in(0,\pi)\)
\begin{equation}
\lim_{n\to\infty}\left(
T_n(x_0)-\frac1\pi\int_0^\delta
\bigl(f(x_0-t)+f(x_0+t)\bigr)D_n(t)\,dt
\right)=0.
\label{eq-fourier-localization}
\end{equation}
Следовательно, сходимость ряда Фурье в точке \(x_0\) и значение его суммы определяются поведением функции в сколь угодно малой окрестности этой точки.

\end{theorem}

\begin{proof}

По лемме \ref{lem-fourier-symmetric-integral} разность в \eqref{eq-fourier-localization} равна
\[
\begin{aligned}
\frac1\pi\int_\delta^\pi\bigl(f(x_0-t)+f(x_0+t)\bigr)D_n(t)\,dt
&=\frac1{2\pi}\int_\delta^\pi h(t)\sin\bigl((n+\tfrac12)t\bigr)\,dt,\\
h(t)&=\frac{f(x_0-t)+f(x_0+t)}{\sin(t/2)}.
\end{aligned}
\]
На \([\delta,\pi]\) знаменатель отделён от нуля: \(\sin(t/2)\ge\sin(\delta/2)>0\). Поэтому
\[
\int_\delta^\pi|h(t)|\,dt
\le\frac1{\sin(\delta/2)}\int_\delta^\pi
\bigl(|f(x_0-t)|+|f(x_0+t)|\bigr)\,dt<+\infty.
\]
Лемма \ref{lem-riemann-oscillatory} при \(\omega=n+1/2\) показывает, что осциллирующий интеграл стремится к нулю. При наличии особых точек применяем её отдельно на соответствующих участках. Получаем \eqref{eq-fourier-localization}.

\end{proof}

\begin{cor}[Совпадение функций в окрестности]\label{cor-fourier-local-agreement}

Пусть \(f\) и \(g\) --- \(2\pi\)-периодические абсолютно интегрируемые функции, совпадающие при \(0<|x-x_0|<\delta\) для некоторого \(\delta>0\). Тогда \(T_n^f(x_0)-T_n^g(x_0)\to0\). Их ряды Фурье сходятся в \(x_0\) одновременно и при сходимости имеют одну и ту же сумму.

\end{cor}

Применяем теорему \ref{thm-fourier-localization} к \(f-g\), выбрав меньшую \(\delta<\pi\). Локальный интеграл равен нулю, откуда следует утверждение.
